\documentclass{beamer}
\usepackage{graphicx}
\usepackage[utf8]{inputenc}
\graphicspath{{Imagenes/}}
%\usetheme{Antibes}
%usetheme{AnnArbor}
%\usetheme{Berkeley}
%\usetheme{CambridgeUS}
%\usetheme{Goettingen}
\usetheme{Bergen}


%Esto unicamente se utiliza para el tema Bergen
%\def\insertauthorindicator{¿Quién?}
%\def\insertdateindicator{Fecha}
%\author{Jorge S. Martínez Villafa}
%\date{\today}
\title{ Taller de Herramientas Computacionaes}
\title{ Taller de Herramientas Computacionaes}
\author{Jorge S. Martínez Villafan}
\date{Enero 22, 2019}

\begin{document}]
	\maketitle
\begin{frame}
\frametitle{Mi primera presentación en LateX}
\end{frame}
%\includegraphics[scale=0.50]{1.png}

\begin{frame}
\frametitle {
Esta es mi segunda diapositiva}
\end{frame}
\begin{frame}[fragile]
\begin{verbatim}
#!/usr/bin/python2.7
# -*- coding: utf-8 -*-
"""
Jorge S Martínez Villafan, 418003839
Herramientas computacionales
Lo que nos explico el miercoles
"""
x = 10.5;y = 1.0/3;z = 15.3
H = """
El punto de R3 es:
(x,y,z)=(%.2f,%g,%G)
""" % (x,y,z)
print H

G="""
El punto en R3 es:
(x,y,z)=({laX:.2f}.{laY:g},{laZ:G})
""".format(laX=x,laY=y,laZ=z)

print G

import math as m
from math import sqrt
from math import sqrt as s
from math import *
x=16
x=input("Cuál es el valor al que le quieres\n" + "calcular la raíz")
print "La raiz cuadrada de %.2f es %f" % (x,m.sqrt(x))
print sqrt(16.5)
print s(16.5)

\end{verbatim}
\end{frame}
\end{document}
